\chapter{Bulk billing and medicare}

\section*{Definition and Overview}
Bulk billing and medicare relates to the organisation, delivery, or accessibility of health and support services. Effective healthcare requires understanding what services exist, when to use them, and how to navigate complex systems. This topic covers the structures, pathways, and resources available to support health and wellbeing.

\section*{Epidemiology}
Health service utilisation is rising with ageing populations, increasing chronic disease, and technological advances. Inequities in access -- geographic, financial, cultural, and informational -- remain a major barrier to care. Emergency departments are frequently overused for non-urgent problems, while preventive and primary care services are underutilised by those who need them most. Health literacy is a strong predictor of appropriate service use.

\section*{Aetiology and Pathophysiology}
\begin{itemize}
    \item \textbf{System complexity:} multiple providers, funding streams, and entry points create confusion
    \item \textbf{Access barriers:} cost, distance, waiting times, language, and cultural factors
    \item \textbf{Information gaps:} people do not know what is available or how to access it
    \item \textbf{Fragmentation:} poor communication between providers and services
\end{itemize}

\section*{Clinical Features}
\begin{itemize}
    \item Uncertainty about which service to use for a given problem
    \item Difficulty obtaining appointments, transport, or affordable care
    \item Feeling that needs are not being met by current care arrangements
    \item Navigating multiple providers and repeated assessments
\end{itemize}

\section*{Differential Diagnosis}
When access to care is difficult, consider the barriers:
\begin{itemize}
    \item Structural: no local service, long waiting list
    \item Financial: out-of-pocket cost, no insurance/concession
    \item Informational: not knowing what is available or how to ask
    \item Cultural: language, trust, or stigma barriers
    \item Personal: transport, mobility, or caring responsibilities
\end{itemize}

\section*{When to Seek Medical Care}
\begin{itemize}
    \item \textbf{Call 000 for emergencies} -- chest pain, severe bleeding, difficulty breathing, collapse, or stroke symptoms
    \item Contact a GP for non-urgent health concerns
    \item Use after-hours services, nurse helplines, or pharmacists for less urgent needs
    \item Ask a community worker or social worker for help navigating services
\end{itemize}

\section*{Diagnosis and Investigations}
\begin{itemize}
    \item \textbf{Needs assessment:} identify what help is needed and roughly which service fits
    \item \textbf{Eligibility check:} for programs, benefits, and concessions
    \item \textbf{Service mapping:} identify available local providers and programs
    \item \textbf{Documentation:} keep a list of current providers, medicines, and care plans
\end{itemize}

\section*{Management}
\begin{itemize}
    \item \textbf{Start with a GP:} the usual entry point and coordinator of ongoing care
    \item \textbf{Use appropriate services:} nurse-led clinics, pharmacist advice, after-hours services for non-urgent needs
    \item \textbf{Explore community services:} peer support groups, carer support, and community health programs
    \item \textbf{Plan ahead:} advance care plans, known contacts, and documented preferences
    \item \textbf{Use emergency services appropriately:} reserve 000 for genuine emergencies
\end{itemize}

\section*{Complications}
\begin{itemize}
    \item Delayed diagnosis and treatment due to access barriers
    \item Inappropriate use of emergency services for non-urgent problems
    \item Fragments of care across multiple uncoordinated providers
    \item Missed preventive care and screening opportunities
    \item Financial hardship from healthcare costs
    \item Disparity in health outcomes for disadvantaged groups
\end{itemize}

\section*{Prognosis and Outcomes}
Effective navigation of health services improves health outcomes, reduces unnecessary costs, and enhances patient experience. A strong relationship with a regular GP is associated with better chronic disease management and reduced hospitalisation. Health literacy and advance planning empower individuals to make informed decisions and receive care aligned with their values.

\section*{Prevention and Self-Care}
\begin{itemize}
    \item Find and keep a regular GP
    \item Keep a list of medicines, allergies, and key contacts
    \item Ask about bulk-billing and concession options early
    \item Build support networks before a crisis
    \item Develop a personal health record and advance care plan
    \item Improve health literacy: ask questions and seek reliable information
\end{itemize}
